\par\Needspace{6\baselineskip}\medskip
\begin{center}\hypertarget{native-181-article135}{}\pdfbookmark[2]{135.}{article-135}135.\end{center}
Disquisitiones praecedentes eum in finem allatae sunt, ut principia, quibus
methodus nostra innititur, verusque eius quasi nervus eo clarius perspiciantur:
tractatio ipsa autem method\'um in forma prorsus diversa exhibebit, quam post
applicationes frequentissimas tamquam commodissimam inter plures alias a
nobis tentatas commendare possumus. Quum in determinanda orbita incognita
e tribus observationibus totum negotium semper ad aliquot hypotheses aut
potius approximationes successivas reducatur, pro lucro eximio habendum erit,
si calculum ita adornare successerit, ut iam ab initio operationes quam plurimas,
quae non a $P$ et $Q$ sed unice a combinatione quantitatum cognitarum
pendeant, ab ipsis hypothesibus separare liceat. Tunc manifesto has operationes
praeliminares, singulis hypothesibus communes, semel tantum exsequi
oportet, hypothesesque ipsae ad operationes quam paucissimas reducuntur.
Perinde maximi momenti erit, si in singulis hypothesibus usque ad ipsa elementa
progredi haud opus fuerit, horumque computum usque ad hypothesin
\hypertarget{native-182}{}postremam reservare liceat. Utroque respectu methodus nostra, quam exponere
iam aggredimur, nihil desiderandum relinquere videtur.
